% No preamble. Source bridge for Notebook18. Unique prefix ib-.
\section{Using the published localization and integration theorem}
\label{ib-section}

There is an existing specialized integration route in the literature.
Ba\l aban's \emph{Renormalization group approach to lattice gauge
field theories II}, Commun. Math. Phys. \textbf{116} (1988),
Lemma 2, pp.~10--11, equations (1.41)--(1.43), localizes the
complete fluctuation exponent, including the difference of inherited
effective interactions. Lemma 3, p.~20, equation (2.38), and the
logarithmic estimate (2.41), p.~21, integrate these overlapping
terms. These results can be used in their published setting;
another abstract integration theorem is not a prerequisite for
using them.

\subsection{The actual coefficient and field domains}

The original domain is equation (1.34), p.~9:
\[
 \mathcal U^c_{k+1}
   (Y,(1+\beta)\alpha_0,(1+\beta)\alpha_1,\alpha_0)
       \;\times\;
 \{B:|B|<\epsilon_1/g\text{ on }Y\}.
\]
The first factor is an orbit domain from CMP109, pp.~262--263,
equations (1.10)--(1.16). Its restrictions include reconstructed
backgrounds on subregions. It cannot be replaced solely by an
unverified global supremum condition. The operator and minimizer
walk deformations in CMP116, pp.~3--6, are also inputs.

To distinguish two scales, let $M_{\rm loc}$ denote the fixed
localization cube size in that paper; it is different from the
refinement $M=L^k$ in our flat covariance theorems. The local
exponent is organized as
\[
 P^{(k)}+\text{inherited-action difference}=\sum_Y V_Y,
 \qquad
 V_Y=\tfrac12\langle B,Q_Y(U,J,B)B\rangle+V_Y''(U,J,B).
\]
The displayed source bounds are (1.43) for the quadratic kernel
and (1.36) for the remainder. Their localization weight is a
normalized minimum tree length meeting the $M_{\rm loc}$ cubes
of $Y$. It is stronger than a diameter bound and comes with
the source's specific geometric summations. The source exponent
contains the Haar, constraint Jacobian and gauge contributions;
it is not only the bare Wilson cubic remainder.

The original article was inspected through its extracted text.
Some displayed numerical coefficients are lost in that extraction;
the numbered source bounds are therefore located here rather than
reconstructed by guesswork. The matching deposited formulas are
Raw I:30590--30597 and Raw II:3190--3208. The latter attributes
(1.43) to CMP109, but that formula in this argument is CMP116
(1.43), as also recorded at Raw II:11178--11188.

\subsection{The logarithmic-radius transfer}

Fix all geometric and analytic parameters before choosing the
coupling ceiling. For $R(g)=c(\log g^{-2})^p$, $p>1/2$,
\[
 gR(g)<\epsilon_1/2,
 \qquad gR(g)^3\longrightarrow0,
 \qquad e^{-c_1R(g)^2}=o(g^q)
\]
for every fixed $q>0$. Thus cutoff points in the actual source
variable lie inside its fixed physical-field analytic radius.
For localized coefficients holomorphic on that radius and
vanishing at the indicated orders, radial Schwarz estimates give
\[
 \|Q_Y(gB)\|\le a_Y gR,
 \qquad |V_Y''(gB)|\le a_Y' gR.
\]
Fixed inverse powers of the analytic radius are included in
$a_Y,a_Y'$. On the cutoff this bounds $V_Y$ by $gR^3$ times
a local coefficient and a polynomial support factor. A reserved
part of the tree-decay exponent pays for that polynomial.
Raw II:11178--11195 uses this distinction between cubic and
merely vanishing contributions.

Changing the real cutoff nevertheless changes the integral.
In the original integration proof the Gaussian bad-bond gain
must be rechecked at the new threshold in (2.22), (2.26) and
(2.31). Its numerical scale is $e^{-cR(g)^2}$, which pays all
fixed geometric entropy below a sufficiently small coupling
ceiling. This numerical fact alone does not verify preservation
of the source's localized operator representations. The
boundary-specific domain CMP119 (3.43) is not a replacement
for the whole-lattice interior domain (1.34).

\subsection{Why the tunable density estimate matters}

The first overlap theorem above has a sufficient interaction norm
containing $K_G^{|Y|}$. If a source localization cube is one cell,
its real coordinate dimension is $D=O(M_{\rm loc}^4)$, and the
original $\log K_G$ has that dimension cost. Taking current-unit
cells instead produces the same loss through their number in
$Y$. The published coefficient decay does not automatically pay
this fixed cost: its parameter hierarchy may choose $M_{\rm loc}$
large after the decay exponent has been fixed. Decreasing the
coupling cannot cure an unfavorable exponential support exponent.

The tunable real-baseline comparison proved in this notebook
reduces that cost to $O(\theta D)$ per source cell. Choose
$\theta$ of order $D^{-1}$, with the source cubes as cells.
Its Cauchy derivative constant then grows with $\theta^{-1}$,
but is fixed before the coupling is decreased. This removes
the automatic $M_{\rm loc}^4$ loss in the support exponent.
The strict version also requires the displayed precision
perturbation smallness at that chosen $\theta$.

For example, if a source coefficient bound has sufficient spare
tree decay to pay the resulting fixed support weight and its
geometric summation, the Schwarz estimate above implies
$\epsilon_\theta(g)\le C_\theta gR(g)^3$.
The interaction and Gaussian-tail smallness conditions then hold
together below a fixed coupling ceiling. This is a precise
conversion criterion; the required source decay and precision
domain must still be matched. Alternatively one can retain
CMP116's own integration theorem, whose controllably small
quadratic and determinant errors serve the same purpose.

\subsection{The immediate model-specific proof}

The flat current-scale covariance is supplied by Notebook 17.
Notebook 19 constructs its nonflat physical-background version,
the exact relative-field pivot chart, and the local measure
factors on the homogeneous small-global-potential branch.
Notebook 20 verifies the full local coefficient and precision
match for the bare first-step Wilson small-field integral.
To invoke the original CMP116 theorem on its larger domains,
one must still show that the background, independent current
and every restricted reconstruction lie in its orbit domain,
and match the source's deformed operators.
If a cutoff or inherited carrier is changed, its exact transfer
through that proof must be established.

The retained cutoff coordinates in Notebook 19 are preserved
exactly. Additional indicators retain their own pullbacks.
The direct first-step proof does not establish the independent
current and mixed-history domains, transport of support to the
next coarse background, or the later quantum interactions.
Those are further model statements needed to close the induction.

